\documentclass[12pt,a4paper]{article}

% ---------- Codificación e idioma ----------
\usepackage[utf8]{inputenc}
\usepackage[T1]{fontenc}
\usepackage[spanish,es-noshorthands]{babel}
\usepackage{lmodern}

% ---------- Paquetes útiles ----------
\usepackage[margin=2.5cm]{geometry}
\usepackage{graphicx}
\usepackage{xcolor}
\usepackage{booktabs}
\usepackage{array}
\usepackage{enumitem}
\usepackage{parskip}
\usepackage{verbatim}
\usepackage{tikz}
\usetikzlibrary{arrows.meta,positioning}
\usepackage[hidelinks]{hyperref}
\usepackage{xurl}   % permite cortar las URLs largas para que no se salgan del margen

% ---------- Referencias con sangria francesa (APA) ----------
\newenvironment{referencias}
  {\begin{list}{}{%
     \setlength{\leftmargin}{1.6em}\setlength{\itemindent}{-1.6em}%
     \setlength{\labelwidth}{0pt}\setlength{\labelsep}{0pt}%
     \setlength{\itemsep}{0.7em}\setlength{\parsep}{0pt}}}
  {\end{list}}

% ---------- Colores (solo para los diagramas) ----------
\definecolor{azulUVG}{RGB}{30,90,150}
\definecolor{verde}{RGB}{46,139,87}
\definecolor{grisclaro}{RGB}{240,243,246}
\definecolor{gristexto}{RGB}{90,100,110}

% Secciones sin numerar para un estilo menos formal
\setcounter{secnumdepth}{0}
\pagestyle{plain}

\begin{document}

% =====================================================================
%  CARÁTULA
% =====================================================================
\begin{titlepage}
\begin{center}

\textbf{UNIVERSIDAD DEL VALLE DE GUATEMALA}\\
Facultad de Ingeniería\\[2cm]

\IfFileExists{logo uvg.png}{%
  \includegraphics[width=0.35\textwidth]{logo uvg.png}\\[2.5cm]
}{%
  \vspace*{2.5cm}
}

{\LARGE\textbf{Publicación de un paquete de Python en TestPyPI}}\\[1.5cm]

Diego Linares\\[0.2cm]
Andy Fuentes\\[0.2cm]
Christian Echeverria\\[0.2cm]
Diederich Solis\\[1.5cm]

Machine Learning Engineering\\[2cm]

Septiembre 1 de 2026

\end{center}
\end{titlepage}

\begin{center}
\textbf{Paquete publicado:}\\
\url{https://test.pypi.org/project/act3-pipeline-mlops/}
\end{center}

\vspace{1em}

\section{Qué decidimos publicar}

Lo primero que tuvimos que aclarar fue \emph{qué} se sube a PyPI, porque nuestra primera idea fue
publicar el trabajo de la Actividad 4, que está dockerizado. No se puede, y por una razón de
fondo: PyPI distribuye \textbf{paquetes de Python}; las imágenes de contenedor van a un registro
de imágenes como Docker Hub. Son dos cosas distintas.

Lo que sí es publicable de la Actividad 4 es la \textbf{librería que vive adentro de sus
imágenes}: el paquete \texttt{act3\_pipeline}, que armamos en la Actividad 1 y que contiene el
pipeline completo con la calibración de hiperparámetros de la Actividad 3.

Y ahí encontramos la mejor justificación para hacerlo. Hoy los dos \texttt{Dockerfile} de la
Actividad 4 instalan la librería copiando un archivo a mano dentro del repositorio:

\begin{verbatim}
COPY libreria/act3_pipeline-0.1.0-py3-none-any.whl /tmp/
RUN pip install --no-cache-dir /tmp/act3_pipeline-0.1.0-py3-none-any.whl
\end{verbatim}

Con el paquete publicado, eso se vuelve una dependencia normal, igual que cualquier otra:

\begin{verbatim}
RUN pip install act3-pipeline-mlops
\end{verbatim}

Se acabó el \texttt{.whl} versionado a mano dentro del repositorio.

\section{La estructura del proyecto}

Seguimos la estructura del tutorial oficial (Python Packaging Authority, s.f.-a): disposición
\texttt{src/} y toda la configuración en un solo \texttt{pyproject.toml}, sin
\texttt{setup.py} ni \texttt{setup.cfg}.

\begin{verbatim}
Tarea4/
|-- pyproject.toml          # toda la configuracion del paquete
|-- README.md               # se muestra como descripcion en TestPyPI
|-- LICENSE                 # MIT
|-- .env                    # el token (NO se sube al repositorio)
|-- .gitignore
|-- src/
|   `-- act3_pipeline/      # el paquete
|       |-- __init__.py
|       |-- data.py         transformers.py    pipeline.py
|       |-- tuning.py       evaluation.py      cli.py
|       `-- datasets/champions_league_matches.csv
`-- dist/                   # lo que genera la construccion
\end{verbatim}

\subsection{El \texttt{pyproject.toml}}

\begin{verbatim}
[build-system]
requires = ["hatchling >= 1.26"]
build-backend = "hatchling.build"

[project]
name = "act3-pipeline-mlops"
version = "0.1.0"
authors = [
    { name = "Diego Linares", email = "lin221256@uvg.edu.gt" },
    { name = "Andy Fuentes" },
    { name = "Christian Echeverria" },
    { name = "Diederich Solis" },
]
description = "Pipeline de scikit-learn con calibracion de hiperparametros ..."
readme = "README.md"
requires-python = ">=3.10"
license = "MIT"
license-files = ["LICENSE"]
classifiers = [
    "Programming Language :: Python :: 3",
    "Operating System :: OS Independent",
    "Intended Audience :: Education",
    "Topic :: Scientific/Engineering :: Artificial Intelligence",
]
dependencies = [
    "scikit-learn>=1.3", "pandas>=2.0", "numpy>=1.24", "scipy>=1.10",
]

[project.scripts]
act3-demo = "act3_pipeline.cli:main"

# El nombre de la distribucion no coincide con el del paquete,
# asi que hay que decirle a hatchling donde buscarlo.
[tool.hatch.build.targets.wheel]
packages = ["src/act3_pipeline"]
\end{verbatim}

Tres detalles que nos costaron y vale la pena anotar:

\begin{itemize}[leftmargin=1.2em]
\item \textbf{El nombre de la distribución no es el nombre del paquete.} Se instala como
\texttt{act3-pipeline-mlops} pero se importa como \texttt{act3\_pipeline}. Al no coincidir, hubo
que declarar \texttt{[tool.hatch.build.targets.wheel]} para que el constructor encontrara el
código.
\item \textbf{El nombre tiene que ser único en todo el índice.} Le agregamos el sufijo
\texttt{-mlops} porque \texttt{act3-pipeline} a secas se podía chocar con otro proyecto.
\item \textbf{\texttt{dependencies} no es lo mismo que \texttt{requirements.txt}.} Aquí van rangos
abiertos (\texttt{>=}) porque quien instale el paquete puede tener otras librerías que también
necesitan resolverse; las versiones exactas se fijan en la aplicación, no en la librería.
\end{itemize}

\section{Construcción}

\begin{verbatim}
python -m build
\end{verbatim}

Genera los dos formatos que espera PyPI:

\begin{center}
\begin{tabular}{|l|l|p{6.2cm}|}
\hline
\textbf{Archivo} & \textbf{Tamaño} & \textbf{Qué es} \\
\hline
\texttt{act3\_pipeline\_mlops-0.1.0-py3-none-any.whl} & 24 KB & Distribución construida; es la que
\texttt{pip} instala directo \\
\hline
\texttt{act3\_pipeline\_mlops-0.1.0.tar.gz} & 20 KB & Código fuente, por si hay que reconstruir \\
\hline
\end{tabular}
\end{center}

Verificamos que el dataset viajara adentro, porque de nada sirve el paquete si hay que descargar
los datos aparte:

\begin{verbatim}
act3_pipeline/__init__.py
act3_pipeline/cli.py           data.py       evaluation.py
act3_pipeline/pipeline.py      tuning.py     transformers.py
act3_pipeline/datasets/champions_league_matches.csv   <-- los datos van adentro
act3_pipeline_mlops-0.1.0.dist-info/METADATA
act3_pipeline_mlops-0.1.0.dist-info/entry_points.txt  <-- el comando act3-demo
act3_pipeline_mlops-0.1.0.dist-info/licenses/LICENSE
\end{verbatim}

Antes de subir nada, validamos la metadata:

\begin{verbatim}
python -m twine check dist/*

Checking dist\act3_pipeline_mlops-0.1.0-py3-none-any.whl: PASSED
Checking dist\act3_pipeline_mlops-0.1.0.tar.gz: PASSED
\end{verbatim}

\section{Publicación en TestPyPI}

Siguiendo la guía de TestPyPI (Python Packaging Authority, s.f.-b), creamos la cuenta, generamos
un token de API y subimos con \texttt{twine}.

\subsection{El manejo del token}

Un token de PyPI permite publicar en nombre de uno, así que no puede quedar escrito en un comando
ni en el historial de la terminal. Lo guardamos en un archivo \texttt{.env} excluido por
\texttt{.gitignore}:

\begin{verbatim}
TWINE_USERNAME=__token__
TWINE_PASSWORD=pypi-AgENdGVzdC5weXBpLm9yZw...
\end{verbatim}

El usuario es literalmente \texttt{\_\_token\_\_}, no el correo. Un script de PowerShell carga
esas variables, construye, valida y sube, de modo que el token nunca aparece en pantalla:

\begin{verbatim}
.\subir_a_testpypi.ps1

Token cargado desde .env (usuario: __token__)
[1/3] Construyendo el paquete...
[2/3] Validando la metadata...
[3/3] Subiendo a TestPyPI...
\end{verbatim}

\IfFileExists{capturas/terminal.png}{%
\begin{figure}[ht]
\centering
\includegraphics[width=\textwidth]{capturas/terminal.png}
\caption{La construcción, la validación y la subida en una sola corrida. Abajo se ve la
instalación posterior desde TestPyPI descargando el \emph{wheel} publicado.}
\end{figure}
}{}

\IfFileExists{capturas/testpypi.png}{%
\begin{figure}[ht]
\centering
\includegraphics[width=\textwidth]{capturas/testpypi.png}
\caption{El paquete ya publicado en TestPyPI, con la descripción tomada del \texttt{README.md},
los enlaces al repositorio y la versión 0.1.0.}
\end{figure}
}{}

\section{Verificación: instalarlo desde el índice}

Publicar no prueba nada; había que instalarlo desde TestPyPI en limpio. Aquí está el error más
común de esta guía, y lo cometimos:

\begin{verbatim}
pip install --index-url https://test.pypi.org/simple/ act3-pipeline-mlops
\end{verbatim}

Eso falla. TestPyPI es un índice \emph{separado} y no tiene scikit-learn ni pandas, así que
\texttt{pip} no puede resolver las dependencias. Hay que decirle que las busque en el PyPI real:

\begin{verbatim}
pip install --index-url https://test.pypi.org/simple/ \
            --extra-index-url https://pypi.org/simple/ \
            act3-pipeline-mlops

Downloading https://test-files.pythonhosted.org/packages/.../
             act3_pipeline_mlops-0.1.0-py3-none-any.whl (24 kB)
Successfully installed act3-pipeline-mlops-0.1.0
\end{verbatim}

Y comprobamos que el paquete instalado desde el índice sirve de verdad:

\begin{verbatim}
>>> import act3_pipeline
distribucion instalada: 0.1.0
se importa como       : act3_pipeline 0.1.0
dataset incluido      : 115 entrenamiento / 29 prueba
pipeline              : ['preprocesamiento', 'seleccion', 'modelo']
\end{verbatim}

El comando de consola también quedó instalado y corre el pipeline completo:

\begin{verbatim}
act3-demo --busqueda aleatoria

RESULTADO - accuracy=0.690  f1_macro=0.631
[OK] Pipeline calibrado y evaluado correctamente en este equipo.
\end{verbatim}

\section{Reflexión: cómo empaquetar un proyecto de Machine Learning}

La pregunta de la tarea es qué conviene más entre Docker, \emph{pickles} y paquetes. Después de
haber usado los tres en el curso, nuestra conclusión es que \textbf{la pregunta está mal
planteada}: no compiten entre sí porque no empaquetan lo mismo.

\begin{center}
\begin{tabular}{|l|p{4.2cm}|p{5.2cm}|}
\hline
\textbf{Herramienta} & \textbf{Qué empaqueta} & \textbf{Dónde la usamos} \\
\hline
Paquete (\emph{wheel}) & El \textbf{código} y sus dependencias declaradas & Actividad 1, y hoy en
TestPyPI \\
\hline
Pickle / joblib & El \textbf{modelo entrenado}: los coeficientes aprendidos & El
\texttt{modelo.joblib} de la Actividad 4 \\
\hline
Docker & El \textbf{ambiente completo}: sistema operativo, intérprete, librerías y código & Las
tres imágenes de la Actividad 4 \\
\hline
\end{tabular}
\end{center}

\subsection{Los tres, en capas}

Lo que nos funcionó es usarlos apilados, cada uno en su papel:

\begin{enumerate}[leftmargin=1.4em]
\item El \textbf{código} se publica como paquete, con sus dependencias declaradas y una versión.
\item El paquete se \textbf{instala dentro de la imagen}, que además fija la versión del
intérprete y del sistema operativo.
\item El \textbf{modelo entrenado vive fuera de la imagen}, en un volumen o en un registro de
modelos, porque es estado y no código.
\end{enumerate}

Ese tercer punto lo aprendimos por decisión propia en la Actividad 4: nos pareció mala idea
hornear el modelo dentro de la imagen porque reentrenar habría obligado a reconstruir y
redistribuir la imagen completa.

\subsection{Por qué el pickle es el eslabón peligroso}

De las tres piezas, la que más problemas nos dio fue el modelo serializado, y por una razón que
solo se entiende cuando pasa: \textbf{un pickle no guarda el código de sus clases, guarda la ruta
para importarlas}. Nuestro pipeline contiene dos transformadores propios,
\texttt{PorcentajeATexto} y \texttt{RatioATexto}. Cuando la API carga el
\texttt{modelo.joblib}, Python busca esas clases en el módulo \texttt{act3\_pipeline}; si no las
encuentra, el archivo es inservible.

Por eso las dos imágenes de la Actividad 4 instalan el mismo paquete aunque la API no entrene
nada, y por eso ambas fijan las mismas versiones exactas. Un modelo serializado sin el paquete
que lo acompaña es una bomba de tiempo: funciona hoy en la máquina de quien lo entrenó y falla
mañana en cualquier otra.

Se suma que \emph{pickle} ejecuta código al deserializar, así que cargar uno de origen
desconocido es un riesgo de seguridad, y que está atado a la versión de la librería con la que se
creó.

\subsection{Nuestra recomendación}

\begin{itemize}[leftmargin=1.2em]
\item \textbf{Si el código se reutiliza entre proyectos} —un pipeline, transformadores, utilidades—
va como \textbf{paquete}. Es lo que hicimos hoy, y elimina el \texttt{.whl} copiado a mano.
\item \textbf{Si hay que reproducir un experimento o servir un modelo}, va en \textbf{Docker}. Es
lo único que garantiza el mismo resultado en Windows y en macOS, algo que sufrimos en la
Actividad 3 antes de contenerizar.
\item \textbf{El modelo entrenado} no va en el repositorio ni horneado en la imagen: va a un
volumen o, mejor, a un \textbf{registro de modelos} con su versión y sus métricas. Al llevar el
pipeline a Databricks vimos cómo MLflow hace exactamente eso, y es bastante más ordenado que
nuestro volumen con un \texttt{metadatos.json} al lado.
\end{itemize}

Si tuviéramos que quedarnos con una sola idea: \textbf{empaquetar el código es fácil y ya está
resuelto}; lo difícil de un proyecto de Machine Learning es empaquetar el modelo y los datos, que
son las dos piezas que cambian solas y que ninguna de estas tres herramientas versiona bien por
sí sola.

% =====================================================================
\section{Referencias}

\begin{referencias}

\item Python Packaging Authority. (s.f.-a). \emph{Packaging Python projects}. Python Packaging
User Guide. \url{https://packaging.python.org/en/latest/tutorials/packaging-projects/}

\item Python Packaging Authority. (s.f.-b). \emph{Using TestPyPI}. Python Packaging User Guide.
\url{https://packaging.python.org/en/latest/guides/using-testpypi/}

\item Python Packaging Authority. (s.f.-c). \emph{Writing your pyproject.toml}. Python Packaging
User Guide. \url{https://packaging.python.org/en/latest/guides/writing-pyproject-toml/}

\item Python Software Foundation. (s.f.). \emph{pickle — Python object serialization}. Python
documentation. \url{https://docs.python.org/3/library/pickle.html}

\item Sculley, D., Holt, G., Golovin, D., Davydov, E., Phillips, T., Ebner, D., Chaudhary, V.,
Young, M., Crespo, J.-F., \& Dennison, D. (2015). Hidden technical debt in machine learning
systems. \emph{Advances in Neural Information Processing Systems}, \emph{28}.

\end{referencias}

\end{document}
